\renewcommand{\abstractname}{ABSTRACT}
\begin{abstract} 
\addcontentsline{toc}{section}{ABSTRACT} 
\begin{center}
\textbf{THE LABOR MARKET IN TURKEY: EMPRICAL INVESTIGATION OF POLARIZATION}
\par\end{center}

\begin{center}
\textbf{Evren Gülser}
\par\end{center}

\begin{center}
\textbf{August, 2023}
\par\end{center}

There is a growing body of literature in the field of labor economics
dedicated to investigating the structural changes in employment over
time. The supply-side perspective relies on the factors such as the
aging of the labor force, increasing female participation, and other
compositional changes influenced by educational upgrading. On the
demand side, the interplay of technology and global trade is the fundamental
explanation. In the early 2000's, a phenomenon called \textquotedbl job
polarization\textquotedbl{} emerged, signifying the pronounced employment
growth at the tails of the skill distribution to encapsulate and define
the general characteristics of employment shifts. Although a consensus
exists regarding the significance of technological advancement as
a primary factor, its impact on labor demand has a level of complexity
that attracts researchers to scrutinize it by employing datasets of
various labor markets. The need for evidence coming from different
labor markets is becoming critical. Thus, this study aims to contribute
to the field by presenting evidence from Turkey. It does so by conducting
an in-depth analysis of the dynamics within Turkey's Labor Market,
utilizing the Labor Force Statistics series provided by TurkStat.
The study period spans from 2004 to 2022, in which we pool 19 waves
of annual cross-sectional micro datasets. The dataset is divided into
two distinct parts: 2004-2012 and 2012-2022, each displaying distinct
characteristics. Consistent with the literature, new datasets are
generated to facilitate a task-based analysis. One dataset covers
individual assessments of job tasks, coming from workers, while the
other dataset incorporates a temporal dimension to comprehensively
capture the evolving composition of tasks over time.

A nearly monotonic relationship between skill level (proxied by median
wage) and employment growth is identified in the first period. In
contrast, the evidence shows a different pattern closer to a polarized
form, albeit substantial real wage increases potentially contaminate
the second period data, thereby obscuring the effect of market forces.
According to evidence, the changing sectoral composition plays a role
in shaping employment shifts. The findings suggest that wage polarization
occurred between 2004 and 2012, and further indicate job polarization
in the period spanning from 2012 to 2022. Occupation-specific characteristics
are sufficient to explain wage polarization in the initial period.
Notably, measurements of tasks at both the occupation and worker levels
serve as meaningful indicators for wage estimation equations.

\end{abstract}

\textbf{Keywords: }labor market dynamics, job polarization, wage polarization,
elasticity of substitution, task-based analysis, labor market, labor
demand, sectoral composition, skill distribution

\textbf{JEL codes: }J21, J24, J31
\begin{center}
\pagebreak
\par\end{center}
